% SGA 3, Expose VI_B, Loop 1 English reconstruction.
% Scope: Corollary 1.3.1 through the complete proof of
% Corollary 1.3.2, including Lemma 1.3.1.1; Proposition 1.4 is excluded.
% Authority: Polo--Gille Exp6B-13oct24.pdf, component-local pages 4--5
% (combined-reader pages 345--346; original-margin printed pages 323--324).
% OCR and the Reinhold English Markdown are comparison witnesses only.

\medskip
\phantomsection\label{cor:sga3-VIB-1.3.1}%
\noindent\textbf{推论 1.3.1.---}
\emph{设 \(A\) 是剩余域为 \(k\) 的 Artin 局部环,
\(G\) 是平坦的 \(A\)-群. 以下断言等价:}%
\footnote{\label{note:sga3-VIB-8}编者注: 仅假设 \(G\) 在某一点处
在 \(A\) 上平坦是不够的. 例如, 假设 \(A\ne k\),
令 \(H\) 为常 \(A\)-群 \((\mathbb Z/2\mathbb Z)_A\),
\(G\) 为 \(H\) 的一个闭 \(A\)-子群,
其非单位分支约化. 则结构态射 \(G\to\operatorname{Spec}A\)
在单位点 \(e\) 处为局部同构, 但在点 \(g\ne e\) 处不平坦.}
\begin{enumerate}
\item[\textup{\SGAInlineRefTarget{cor:sga3-VIB-1.3.1:item:i}(i)}]
  \emph{\(G\) 在某点处在 \(A\) 上局部拟有限
  (分别地, 非分歧、平滑或平展);}
\item[\textup{\SGAInlineRefTarget{cor:sga3-VIB-1.3.1:item:ii}(ii)}]
  \emph{\(G\) 在 \(A\) 上局部拟有限
  (分别地, 非分歧、平滑或平展).}
\end{enumerate}

\footnote{\label{note:sga3-VIB-9}编者注: 添加了下面一句,
并以 EGA~IV, 2.7.1 代替上述文献 17.7.5,
简化了 \hyperref[lem:sga3-VIB-1.3.1.1]{引理~1.3.1.1} 的证明.}
\emph{证明.}
若 \(G\) 在某点 \(x\) 满足
\textup{\SGARefLink{cor:sga3-VIB-1.3.1:item:i}{(i)}} 的某一条件,
则存在 \(x\) 的开邻域 \(U\), 在 \(A\) 上为有限型.
因此, 结合下面的引理, 只须在 \(H\) 为单位 \(A\)-群的情形
应用 \hyperref[prop:sga3-VIB-1.3]{1.3} 即可.

\medskip
\phantomsection\label{lem:sga3-VIB-1.3.1.1}%
\noindent\textbf{引理 1.3.1.1.---}
\emph{设 \(A\) 为 Artin 局部环, \(G\) 为 \(A\)-群.
若 \(G\) 包含一个在 \(A\) 上为有限型的非空开子集,
则 \(G\) 在 \(A\) 上局部有限型.}

\emph{证明.}
由 \hyperref[lem:sga3-VIB-1.3.0]{引理~1.3.0},
可假设 \(A\) 等于其剩余域 \(k\).
此外, 由 \(\mathrm{fpqc}\) 下降, 可假设 \(k\) 代数闭
(参见 EGA~IV\(_2\), 2.7.1).
令 \(V\) 为 \(G\) 中使 \(G\) 在 \(k\) 上为有限型的点组成的开子集;
由假设, \(V\ne\varnothing\). 由于 \(G\) 是 Jacobson 概形,
\(V\) 包含一个闭点 \(x\); 为证明 \(V=G\),
只须证明 \(G\) 的每个闭点 \(y\) 都属于 \(V\).
对这样的点 \(y\), 平移
\[
  r_y\circ r_x^{-1}
\]
将 \(x\) 映到 \(y\), 因而 \(y\in V\).

\medskip
\phantomsection\label{cor:sga3-VIB-1.3.2}%
\noindent\textbf{推论 1.3.2.---}
\emph{设 \(A\) 为 Artin 局部环, \(u:G\to H\)
为局部有限型的 \(A\)-群之间的态射. 以下断言等价:
\begin{enumerate}
\item[\textup{\SGAInlineRefTarget{cor:sga3-VIB-1.3.2:item:i}(i)}]
  \(u\) 普遍开;
\item[\textup{\SGAInlineRefTarget{cor:sga3-VIB-1.3.2:item:ii}(ii)}]
  \(u\) 开;
\item[\textup{\SGAInlineRefTarget{cor:sga3-VIB-1.3.2:item:iii}(iii)}]
  \(u\) 在 \(G\) 的某点开;
\item[\textup{\SGAInlineRefTarget{cor:sga3-VIB-1.3.2:item:iv}(iv)}]
  \(u\) 诱导的态射 \(u^0:G^0\to H^0\) 支配;
\item[\textup{(iv\('\))}] \(u^0\) 满射;
\item[\textup{\SGAInlineRefTarget{cor:sga3-VIB-1.3.2:item:v}(v)}]
  存在 \(G\) 的连通分支 \(C\), 使得若 \(D\)
  表示 \(H\) 中包含 \(u(C)\) 的连通分支,
  则诱导态射 \(u':C\to D\) 支配.
\end{enumerate}}

\emph{证明.}
蕴含关系
\[
  \textup{(i)}\Longrightarrow\textup{(ii)}
  \Longrightarrow\textup{(iii)}
  \qquad\text{以及}\qquad
  \textup{(iv\('\))}\Longrightarrow\textup{(iv)}
  \Longrightarrow\textup{(v)}
\]
是显然的. 由于 \(G^0\) 在 \(A\) 上为有限型
(第~VI\(_A\)~篇, \SGARefLink{cor:sga3-VIB-2.4}{2.4}),
因而为 Noether 概形, 所以 \(u^0\) 拟紧.
因此, 由 \hyperref[prop:sga3-VIB-1.2]{1.2},
\(u^0(G^0)\) 在 \(H^0\) 中闭, 故
\textup{\SGARefLink{cor:sga3-VIB-1.3.2:item:iv}{(iv)}} 蕴含 \textup{(iv\('\))}.
另一方面, \(G^0\) (分别地, \(C\)) 在 \(G\) 中开
(第~VI\(_A\)~篇, \SGARefLink{cor:sga3-VIB-2.3}{2.3}),
而 \(H^0\) (分别地, \(D\)) 不可约
(第~VI\(_A\)~篇, \SGARefLink{cor:sga3-VIA-2.4.1}{2.4.1}).
因而 \textup{\SGARefLink{cor:sga3-VIB-1.3.2:item:ii}{(ii)}} 蕴含
\textup{\SGARefLink{cor:sga3-VIB-1.3.2:item:iv}{(iv)}}
(分别地, \textup{\SGARefLink{cor:sga3-VIB-1.3.2:item:iii}{(iii)}} 蕴含
\textup{\SGARefLink{cor:sga3-VIB-1.3.2:item:v}{(v)}}).
只须证明 \textup{\SGARefLink{cor:sga3-VIB-1.3.2:item:v}{(v)}} 蕴含
\textup{\SGARefLink{cor:sga3-VIB-1.3.2:item:i}{(i)}}.

赋予 \(G\) (分别地, \(H\)) 的开子集 \(C\) (分别地, \(D\))
以诱导的概形结构, 并用 \(u'\) 表示诱导态射 \(C\to D\).
令 \(k\) 为 \(A\) 的剩余域.%
\footnote{\label{note:sga3-VIB-10}编者注: 以下展开了原文.}
由于基变换
\[
  \operatorname{Spec}k\longrightarrow\operatorname{Spec}A
\]
为普遍同胚, 可假设 \(A=k\).
由假设 \(u'\) 支配; 又由于 \(C\) 在 \(k\) 上为有限型
(第~VI\(_A\)~篇, 2.4.1), 因而为 Noether 概形, 所以 \(u'\) 拟紧.
由 EGA~IV\(_2\), 2.3.7,
\[
  u'\otimes_k\bar{k}
\]
仍然拟紧且支配, 其中 \(\bar{k}\) 表示 \(k\) 的一个代数闭包.
由于 \(C\otimes_k\bar{k}\) 是 \(G\otimes_k\bar{k}\)
的一些连通分支的并, 态射
\[
  u\otimes_k\bar{k}:
  G\otimes_k\bar{k}\longrightarrow H\otimes_k\bar{k}
\]
满足断言 \textup{(v)}.
根据 EGA~IV\(_2\), 2.6.4, 因而化归到 \(A=k\) 为代数闭域的情形.

在此情形中, 还可将 \(u\) 替换为 \(u_{\mathrm{red}}\),
从而化归到 \(H\) 约化的情形.
令 \(\xi\) (分别地, \(\eta\)) 为 \(C\) (分别地, \(D\)) 的泛点.
由于 \(u'\) 拟紧且支配,
\[
  u'(\xi)=\eta
\]
(参见 EGA~IV\(_1\), 1.1.5).
由于 \(H\) 约化, 局部环 \(\mathcal O_{H,\eta}\) 是域,
因而 \(u'\) 在 \(\xi\) 处平坦.%
\footnote{\label{note:sga3-VIB-11}编者注:
原文调用了泛平坦性定理 (EGA~IV\(_2\), 6.9.1), 这里并不需要它.}
由 \SGARefLink{prop:sga3-VIB-1.3}{1.3}, \(u\) 平坦.
此外, 由于 \(u\) 局部有限型, 且 \(H\) 局部 Noether,
\(u\) 局部有限呈示. 因此由 EGA~IV\(_2\), 2.4.6, \(u\) 普遍开.
% Translation statistics: source_lines=152; translated_lines=146; inline_math=128; display_math=6; pending_TODO=0.
% Review: complete component translation; math fragments, reference keys, environments and note counts checked against source; formula prose translated; compilation pending parent integration.
